\chapter{PENDAHULUAN}
\pagenumbering{arabic}

\section{Latar Belakang}

Formalin (larutan formaldehida) merupakan bahan kimia berbahaya yang dilarang keras penggunaannya sebagai bahan tambahan pangan di Indonesia berdasarkan Peraturan Menteri Kesehatan Republik Indonesia No. 033/2012. Kendati demikian, praktik penyalahgunaannya sebagai pengawet ilegal masih kerap ditemukan pada produk pangan olahan hewani siap saji, salah satunya bakso \cite{menlhk_formalin,padang_bakso,sosis_efsa}. Berbagai studi identifikasi di lapangan secara konsisten melaporkan temuan positif formalin pada sampel bakso yang beredar di pasar tradisional maupun pedagang kaki lima, seperti di Kota Padang \cite{padang_bakso}, sepanjang Jalan Soebrantas Kota Pekanbaru \cite{pekanbaru_bakso}, hingga di area wisata dan kantin sekolah di Yogyakarta \cite{yogyakarta_bakso}. Salah satu penelitian di Pekanbaru bahkan melaporkan kadar formalin pada sampel bakso mencapai 4,506~ppm \cite{pekanbaru_bakso}, padahal Otoritas Keamanan Pangan Eropa (\textit{European Food Safety Authority}/EFSA) menetapkan ambang batas toleransi residu formalin dalam pangan maksimal hanya sebesar 100~ppm \cite{sosis_efsa}. Mengingat formaldehida bersifat karsinogenik golongan 1 (IARC), paparan residu secara berulang meskipun dalam konsentrasi rendah berpotensi terakumulasi dalam jaringan biologis dan memicu kerusakan organ kronis. Kondisi tersebut menegaskan urgensi tersedianya sistem deteksi formalin yang objektif, memiliki sensitivitas tinggi untuk menjangkau ambang batas regulasi pangan, serta mampu memberikan konfirmasi kuantitatif secara langsung di lapangan.

Metode pengujian formalin konvensional saat ini masih didominasi oleh uji kualitatif menggunakan reagen warna kimia (seperti kalium permanganat $\mathrm{KMnO_4}$, pereaksi asam kromotropat, maupun \textit{rapid test kit} komersial) serta metode analitik laboratorium berbasis spektrofotometri UV-Vis dan kromatografi cair kinerja tinggi (HPLC). Meskipun metode laboratorium memiliki akurasi tinggi dan batas deteksi (\textit{limit of detection}/LOD) di bawah 1~mg/L, metode tersebut bersifat destruktif terhadap sampel, memerlukan prosedur derivatisasi kimiawi yang kompleks dan memakan waktu (30--60 menit per sampel), melibatkan reagen kimia toksik (seperti pereaksi Nash atau 2,4-DNPH), serta menggunakan instrumen stasioner yang mahal sehingga tidak fleksibel untuk inspeksi lapangan. Sebaliknya, metode uji cepat kolorimetri di lapangan umumnya hanya menghasilkan indikasi visual biner (ada atau tidaknya perubahan warna) dengan batas deteksi yang relatif kasar (10--50~mg/L), serta sangat rentan terhadap bias interpretasi subjektif dari operator. Kesenjangan antara metode laboratorium yang akurat tetapi lambat dan uji cepat visual yang rentan galat mendorong intensifnya riset pengembangan sensor analitik berbasis elektronik.

Dalam perkembangannya, berbagai sensor elektronik dan biosensor telah dieksplorasi untuk mendeteksi formalin, di antaranya sensor gas berbasis semikonduktor oksida logam (\textit{metal oxide semiconductor}/MOS), biosensor elektrokimia berbasis enzim, dan sensor kolorimetri digital terintegrasi optoelektronik. Sensor gas MOS berbasis $\mathrm{SnO_2}$ atau $\mathrm{In_2O_3}$ bekerja melalui perubahan konduktivitas permukaan akibat desorpsi oksigen saat bereaksi dengan gas formaldehida, menghasilkan batas deteksi pada rentang sub-ppm (0.05--1.0~ppm) \cite{wang2021mos}. Namun, sensor MOS menuntut suhu operasional tinggi (200--400~$^\circ\mathrm{C}$) yang membutuhkan pemanas berdaya besar sehingga tidak ramah baterai portabel, serta memiliki kelemahan mendasar berupa interferensi silang (\textit{cross-sensitivity}) terhadap kelembapan udara dan senyawa volatil bumbu bakso (seperti uap alkohol dan asam asetat). Di sisi lain, biosensor elektrokimia memanfaatkan enzim \textit{formaldehyde dehydrogenase} (FDH) yang menawarkan selektivitas biokimiawi tinggi dengan LOD hingga 0.02~mg/L \cite{singhal2024food}. Kendati demikian, enzim merupakan makromolekul yang sangat labil dengan masa simpan pendek (kurang dari 3 minggu), memerlukan kofaktor $\mathrm{NAD}^+$ yang mahal, sangat sensitif terhadap perubahan suhu dan pH lingkungan, serta rentan mengalami peracunan elektroda (\textit{electrode fouling}) akibat adsorpsi makromolekul lemak dan protein bakso pada permukaan sensor. Sementara itu, sistem sensor optoelektronik dan kolorimetri digital yang mengandalkan analisis warna citra digital \cite{sun2023optical} sangat rentan terdistorsi oleh kekeruhan alami (\textit{turbidity}) dan pigmen alami daging bakso, sehingga tetap membutuhkan preparasi sentrifugasi bertingkat. Keterbatasan-keterbatasan tersebut memicu kebutuhan akan platform sensor alternatif yang mampu beroperasi stabil pada suhu ruang, hemat energi, kebal terhadap warna atau kekeruhan sampel fisik, dan memiliki reseptor sintetis berkestabilan tinggi, yang salah satu solusinya bermuara pada biosensor magnetoresistif berbasis nanokomposit magnetik.

Pemanfaatan nanokomposit magnetik berbasis serat nano (\textit{nanofiber}) menjadi pendekatan inovatif dalam merealisasikan elemen pengenal (\textit{chemo-receptor}) yang selektif dan berdaya tahan tinggi. Nanopartikel \ch{Fe3O4} memiliki keunggulan berupa sifat superparamagnetik, suseptibilitas magnetik tinggi, dan biokompatibilitas yang sangat baik \cite{antarnusa2022,ardiyanti2025}. Melalui teknik elektrospinning, matriks polimer polivinil alkohol (PVA) yang terdispersi nanopartikel \ch{Fe3O4} dapat dibentuk menjadi jejaring nanofiber berpori dengan rasio luas permukaan terhadap volume yang masif \cite{xue2019electrospinning,Turkoglu2024}, sehingga menyediakan densitas situs aktif interaksi analit yang berlimpah. Pada arsitektur nanokomposit ini, asam sitrat berperan ganda sebagai agen penstabil dispersi (\textit{capping agent}) nanopartikel \ch{Fe3O4} sekaligus agen penaut silang hijau (\textit{green crosslinker}) tunggal bagi rantai polimer PVA. Melalui reaksi esterifikasi termal (\textit{thermal curing}) pada suhu $130\,^{\circ}\mathrm{C}$, asam sitrat membentuk jejaring ester kovalen antar-rantai PVA yang memperkokoh integritas mekanik serta ketahanan air (\textit{water resistance}) membran tanpa melepaskan residu toksik \cite{shi2008citric,birck2014citric}. Sementara itu, fungsi pengenalan analit dipercayakan secara khusus kepada \textit{adipic acid dihydrazide} (ADH) yang difungsionalisasikan sebagai reseptor kimiawi (\textit{chemo-receptor}). Berbeda dari glutaraldehida (GA) yang tergolong dialdehida dan rentan memicu reaksi silang kompetitif serta ambiguitas pengenalan analit, ADH memiliki dua gugus terminal hidrazida (\ch{-NH-NH2}) yang bereaksi secara spontan, spesifik, dan kovalen dengan gugus karbonil formaldehida membentuk ikatan hidrazon (\ch{-C=N-NH-}) yang stabil pada suhu ruang \cite{sigma_adh,gantrade_adh}. Pemisahan peran fungsional yang tegas ini---asam sitrat sebagai penaut silang struktural dan ADH sebagai reseptor molekuler---menjamin stabilitas fisik membran sekaligus spesifisitas pengenalan formaldehida yang tinggi.

Interaksi spesifik antara molekul formaldehida dan gugus hidrazida ADH pada permukaan nanofiber memodulasi momen magnetik lokal serta sebaran medan magnetik bocor (\textit{stray field}) dari nanopartikel \ch{Fe3O4}. Deteksi pergeseran sinyal medan magnetik mikro tersebut menuntut transduser magnetik dengan resolusi dan sensitivitas tinggi. Sebagian besar penelitian biosensor magnetik terdahulu memanfaatkan sensor \textit{Giant Magnetoresistance} (GMR) dengan rasio magnetoresistansi $\Delta R/R$ berkisar 10--15\% dan sensitivitas rasiometrik 2--4~mV/V/mT \cite{gilang2026skripsi}. Guna meningkatkan batas resolusi deteksi, penelitian ini menerapkan transduser \textit{Tunneling Magnetoresistance} (TMR) komersial tipe ALT023-10E. Sensor TMR beroperasi berdasarkan fenomena penerobosan terowongan kuantum (\textit{quantum tunneling effect}) melintasi lapisan penghalang isolator magnesium oksida (MgO) berorde nanometer pada struktur \textit{magnetic tunnel junction} (MTJ) \cite{nanomedicine_mtj}. Mekanisme \textit{tunneling} ini menghasilkan rasio magnetoresistansi melampaui 200\% pada suhu ruang serta sensitivitas diferensial 15--25~mV/V/mT, yakni berlipat ganda dibandingkan sensor GMR pada rentang medan magnet rendah $\pm 1.0$~mT \cite{nve_alt023}. Karakteristik sensitivitas tersebut memungkinkan transduser TMR menangkap pergeseran fluks magnetik lokal akibat interaksi formalin pada konsentrasi rendah secara presisi.

Pembacaan sinyal diferensial sensor TMR yang berada pada orde fraksi milivolt membutuhkan rantai instrumentasi presisi tinggi sebelum dikonversi menjadi data digital. Pada penelitian ini, keluaran diferensial sensor TMR dikuatkan menggunakan \textit{instrumentation amplifier} AD623 bertegangan tunggal $+5.0\text{ V}$ dengan tegangan referensi presisi $V_{\text{REF}} = 2.50\text{ V}$ yang dibentuk melalui pembagi tegangan resistor presisi $1.0\text{ k}\Omega$ \cite{adi_ad623}. Rangkaian diperkuat dengan tapis lolos-rendah aktif (\textit{active RC low-pass filter}) berfrekuensi potong $f_c = 10\text{ Hz}$ guna meredam ripple jala-jala listrik 50~Hz serta derau termal lingkungan. Sinyal analog yang telah tersaring didigitasi oleh ADC eksternal ADS1115 16-bit beresolusi tinggi (0.1875~mV/count) \cite{ti_ads1115}, diakuisisi mikrokontroler Arduino Uno secara deterministik berbasis interupsi, dan ditransmisikan secara serial ke modul pemroses utama. Kepastian validitas dan keterulangan (\textit{repeatability}) metrologis dijamin melalui kalibrasi awal sistem menggunakan medan magnet acuan terkontrol dari sepasang kumparan Helmholtz terstandardisasi dan teslameter terkalibrasi \cite{zhu2023,ginisa2015,yudhistira2019}.

Tantangan berikutnya terletak pada karakteristik respons sinyal biosensor terhadap konsentrasi formalin pada matriks pangan nyata (bakso), yang kerap menunjukkan perilaku nonlinearitas, histeresis mikro, dan fluktuasi akibat derau matriks biologis (kandungan garam, lemak, dan protein daging). Mengandalkan regresi linear sederhana semata rentan menimbulkan galat estimasi konsentrasi yang signifikan. Oleh karena itu, diterapkan pemodelan cerdas berbasis ekstraksi tiga fitur sinyal domain waktu utama: tegangan rata-rata kondisi tunak ($V_{\text{mean}}$), pergeseran tegangan respons terhadap baseline ($\Delta V$), dan simpangan baku fluktuasi derau sinyal ($\sigma_V$). Vektor fitur tiga dimensi ($d=3$) ini memberikan representasi kompak yang ideal untuk dipetakan secara terukur ke dalam tiga qubit ($n=3$) pada ruang keadaan Hilbert kuantum berdimensi $2^3 = 8$. Pemodelan dilakukan melalui komparasi komprehensif antara dua paradigma: dua model klasik, yaitu \textit{Support Vector Machine} (SVM) dan \textit{Random Forest} (RF), serta dua model kuantum, yaitu \textit{Quantum Support Vector Classifier} (QSVC) dan \textit{Variational Quantum Classifier} (VQC). Pada paradigma klasik, SVM dengan kernel \textit{Radial Basis Function} (RBF) memetakan batas keputusan nonlinear berdasarkan prinsip \textit{structural risk minimization} \cite{cortes1995support}, sementara Random Forest memanfaatkan agregasi pohon keputusan (\textit{ensemble bagging}) yang tangguh terhadap variansi data sensor \cite{Breiman2001,Tyralis2019}. Pada paradigma kuantum (\textit{Quantum Machine Learning}/QML), QSVC menghitung matriks kemiripan data menggunakan \textit{Quantum Kernel Estimator} berbasis sirkuit penyandian $ZZ\text{FeatureMap}$ \cite{havlicek2019supervised,schuld2019quantum}, sedangkan VQC melatih gerbang rotasi kuantum parametrik (\textit{parameterized quantum circuit}) secara variasional melalui penurunan gradien \cite{schuld2019quantum}. Evaluasi komparatif keempat model menggunakan validasi silang \textit{5-fold cross validation} dengan metrik akurasi, presisi, \textit{recall}, \textit{F1-score}, efisiensi komputasi, dan ketahanan derau (\textit{noise resilience}) menjadi kontribusi analitis penting dalam mengevaluasi potensi komputasi kuantum pada instrumentasi analitik pangan.

Berdasarkan keseluruhan landasan pemikiran di atas, hingga saat ini belum ditemukan penelitian yang secara terpadu mengintegrasikan transduser sensor \textit{Tunneling Magnetoresistance} (TMR) beresolusi tinggi, material fungsional nanofiber \ch{Fe3O4}/PVA-Sitrat-ADH sebagai elemen reseptor kovalen selektif, rantai instrumentasi akuisisi analog-digital AD623-ADS1115 yang stabil, serta komparasi model klasifikasi cerdas dua model klasik (SVM dan RF) versus dua model kuantum (QSVC dan VQC) untuk deteksi formalin pada bakso. Oleh sebab itu, penelitian ini dirancang untuk menjawab tantangan tersebut secara komprehensif, mulai dari fabrikasi material, perancangan mekatronika perangkat keras, kalibrasi metrologis, pengujian analit formalin pada sampel bakso riil, hingga analisis komparasi performa model komputasi cerdas. Hasil penelitian ini diharapkan dapat menjadi fondasi ilmiah yang kokoh dalam melahirkan instrumen biosensor formalin yang portabel, berakurasi tinggi, dan terjangkau guna mendukung penjaminan keamanan pangan nasional.



\section{Rumusan Masalah}
Berdasarkan latar belakang di atas maka permasalahan yang dapat dirumuskan adalah sebagai berikut:
\begin{enumerate}
    \item Bagaimana sensitivitas, batas deteksi (\textit{limit of detection}/LOD), dan pengaruh gugus fungsional hidrazida ADH terhadap karakteristik respons magnetoresistif sensor TMR berbasis nanofiber \ch{Fe3O4}/PVA-Sitrat-ADH dalam mendeteksi variasi konsentrasi formalin?
    \item Bagaimana karakteristik medan magnetik acuan kumparan Helmholtz dan respons linearitas transduser TMR ALT023-10E sebelum dan sesudah integrasi rantai pengkondisi sinyal?
    \item Bagaimana rancangan sistem instrumentasi (AD623-ADS1115-Arduino-Raspberry Pi) menghasilkan kurva kalibrasi yang valid antara tegangan keluaran sensor dan konsentrasi formalin pada bakso?
    \item Bagaimana komparasi performa klasifikasi antara model klasik (\textit{Support Vector Machine} dan \textit{Random Forest}) serta model kuantum (\textit{Quantum Support Vector Classifier} dan \textit{Variational Quantum Classifier}) dalam mendeteksi dan mengklasifikasikan kandungan formalin pada sampel bakso?
\end{enumerate}


\section{Batasan Masalah}
Batasan masalah dalam penelitian ini ditetapkan agar fokus pengujian tetap terarah dan terukur, meliputi:
\begin{enumerate}
    \item Penelitian ini difokuskan khusus pada deteksi analit formalin; bahan tambahan pangan berbahaya lain seperti boraks atau pewarna tekstil tidak termasuk cakupan analisis.
    \item Pengujian sensor diawali menggunakan larutan formalin standar hasil pengenceran bertingkat dengan konsentrasi terkontrol (mg/L) sebagai dasar kurva kalibrasi, sebelum diterapkan pada ekstrak sampel bakso riil.
    \item Material sensor yang digunakan adalah nanofiber \ch{Fe3O4}/PVA-Sitrat-ADH hasil elektrospinning dengan penautan silang termal asam sitrat dan fungsionalisasi gugus reseptor hidrazida ADH.
    \item Mekanisme pengenalan formalin oleh gugus hidrazida ADH pada penelitian ini diperlakukan sebagai hipotesis yang diuji secara empiris melalui respons transduksi magneto-mekanik, bukan mekanisme yang telah dibuktikan secara pasti.
    \item Evaluasi sensor TMR dibatasi pada parameter sensitivitas, linearitas, batas deteksi (LOD), dan kestabilan sinyal; aspek lain seperti ketahanan mekanik jangka panjang tidak menjadi fokus utama.
    \item Pemodelan cerdas dibatasi pada perbandingan dua paradigma: dua model klasik (\textit{Support Vector Machine} ber-kernel RBF dan \textit{Random Forest}) serta dua model kuantum (\textit{Quantum Support Vector Classifier} berbasis $ZZ\text{FeatureMap}$ dan \textit{Variational Quantum Classifier} berbasis sirkuit kuantum parametrik) yang disimulasikan menggunakan pustaka Qiskit pada komputer klasik.
\end{enumerate}


\section{Tujuan Penelitian}
Berdasarkan rumusan masalah di atas, penelitian ini bertujuan untuk:
\begin{enumerate}
    \item Mengevaluasi sensitivitas, batas deteksi (LOD), serta pengaruh fungsionalisasi gugus fungsional hidrazida ADH terhadap karakteristik respons magnetoresistif sensor TMR berbasis nanofiber \ch{Fe3O4}/PVA-Sitrat-ADH dalam mendeteksi variasi konsentrasi formalin.
    \item Mengkarakterisasi medan magnetik acuan kumparan Helmholtz serta respons linearitas transduser TMR ALT023-10E sebelum dan sesudah integrasi rantai pengkondisi sinyal.
    \item Merancang dan memvalidasi sistem instrumentasi akuisisi data (AD623-ADS1115-Arduino-Raspberry Pi) yang menghasilkan kurva kalibrasi valid antara tegangan keluaran sensor dan konsentrasi formalin pada bakso.
    \item Menganalisis dan membandingkan performa klasifikasi (akurasi, presisi, \textit{recall}, \textit{F1-score}, efisiensi komputasi, dan ketahanan derau) antara model klasik (SVM dan RF) serta model kuantum (QSVC dan VQC) dalam mengklasifikasikan kandungan formalin pada sampel bakso.
\end{enumerate}

\section{Manfaat Penelitian}
Penelitian ini akan memberikan beberapa manfaat, di antaranya:
\begin{enumerate}
    \item Memberikan kontribusi teoretis dalam pengembangan ilmu instrumentasi sensor dengan menghadirkan kajian mengenai performa sensor \textit{Tunneling Magnetoresistance} (TMR) berbasis nanofiber \ch{Fe3O4}/PVA-Sitrat-ADH untuk deteksi formalin.
    \item Menawarkan pendekatan baru melalui pemisahan peran fungsional material: asam sitrat sebagai agen penaut silang hijau (\textit{green crosslinker}) ramah lingkungan yang mengeliminasi glutaraldehida toksik, dan \textit{adipic acid dihydrazide} (ADH) sebagai reseptor kovalen spesifik yang selektif mengenali gugus karbonil formaldehida melalui formasi ikatan hidrazon.
    \item Memberikan kontribusi kebaruan ilmiah dalam penerapan komparasi komputasi kuantum (\textit{Quantum Machine Learning}: QSVC dan VQC) terhadap model klasik (SVM dan RF) untuk klasifikasi sinyal sensor pada pengujian keamanan pangan.
    \item Menjadi dasar bagi pengembangan instrumen deteksi formalin yang cerdas, portabel, kuantitatif, dan terjangkau untuk mendukung pengawasan keamanan pangan nasional.
\end{enumerate}

\section{Metode Pengumpulan Data}
Pencapaian tujuan penelitian ditempuh melalui serangkaian metode pengumpulan data sebagai berikut:
\begin{enumerate}
    \item Studi Literatur: Mengumpulkan dan menganalisis informasi dari penelitian terdahulu mengenai deteksi formalin pada pangan, teknologi \textit{Tunneling Magnetoresistance} (TMR), material nanofiber \ch{Fe3O4}/PVA-Sitrat-ADH, serta algoritma pemodelan klasik (SVM dan RF) dan kuantum (QSVC dan VQC).
    \item Observasi dan Studi Pendahuluan: Melakukan observasi terhadap prototipe instrumentasi sensor magnetoresistif yang telah dikembangkan sebelumnya, sebagai titik tolak perancangan rangkaian instrumentasi sensor TMR pada penelitian ini, termasuk teknik fabrikasi nanofiber melalui elektrospinning dengan stabilisasi asam sitrat dan fungsionalisasi ADH.
    \item Eksperimen Laboratorium: Merancang dan menguji instrumen sensor TMR berbasis nanofiber \ch{Fe3O4}/PVA-Sitrat-ADH untuk mendeteksi variasi konsentrasi formalin, diawali karakterisasi medan magnet kumparan Helmholtz hingga pengujian sampel bakso riil.
    \item Analisis Data dan Pemodelan Cerdas: Data hasil eksperimen diolah melalui analisis regresi kalibrasi, ekstraksi fitur sinyal ($V_{\text{mean}}$, $\Delta V$, $\sigma_V$), serta pelatihan dan evaluasi komparatif model klasik (SVM, RF) versus model kuantum (QSVC, VQC).
\end{enumerate}

\newpage
\section{Sistematika Penulisan}
Dalam penulisan penelitian ini, diatur sistematika penulisan dalam tiga bab dengan urutan sebagai berikut:
\\
\\
BAB I: PENDAHULUAN

Bab ini berisi uraian mengenai latar belakang penelitian, perumusan masalah, batasan masalah, tujuan penelitian, manfaat penelitian, serta metode pengumpulan data yang digunakan. Bagian ini menjadi landasan awal untuk menjelaskan pentingnya rancang bangun instrumentasi sensor \textit{Tunneling Magnetoresistance} (TMR) berbasis nanofiber \ch{Fe3O4}/PVA-Sitrat-ADH dengan komparasi model klasik (SVM, RF) dan model kuantum (QSVC, VQC) dalam deteksi formalin pada bakso.
\\
\\
BAB II: TINJAUAN PUSTAKA

Bab ini membahas kajian pustaka dan teori-teori yang relevan dengan penelitian, meliputi permasalahan formalin pada pangan, prinsip kerja sensor \textit{Tunneling Magnetoresistance} (TMR), dasar material pembentuk nanofiber (PVA, Asam Sitrat, ADH, dan nanofiber komposit), komponen instrumentasi perangkat keras, hingga dasar teori pemodelan komparasi model klasik (SVM, RF) dan model kuantum (QSVC, VQC).
\\
\\
BAB III: METODE PENELITIAN DAN PERANCANGAN SISTEM

Bab ini memuat daftar alat dan bahan yang digunakan, tahapan sintesis terstruktur nanofiber komposit melalui elektrospinning dengan parameter optimal, perancangan sistem akuisisi data perangkat keras dan perangkat lunak, serta metodologi ekstraksi fitur dan evaluasi komparatif model klasik (SVM, RF) versus model kuantum (QSVC, VQC).
